\begin{enumerate}
    \item[\textbf{Problem 2.1.}] Justify the assertion made just above that ``nine (9) different graphs, each containing three (3) curves'' are needed to relate force and speed.

    \item[\textbf{Problem 2.2.}] Find and compare the mass density and viscosity of peanut butter, honey, and water.

    \item[\textbf{Problem 2.3.}] What is the constant in eq. (2.12)? How do you know that?

    \item[\textbf{Problem 2.4.}] Apply the basic method to eq. (2.2) for the period of the pendulum.

    \item[\textbf{Problem 2.5.}] Carry out the basic method for eq. (2.13) and show that the mass of a falling body does not affect its speed of descent.

    \item[\textbf{Problem 2.6.}] Carry out the last step of the basic method for eqs. (2.20) using the first mass and show it produces a form that is equivalent to eq. (2.21).

    \item[\textbf{Problem 2.7.}] Write out the Buckingham Pi theorem as a series of steps, analogous to the steps described in the basic method on p. 23.

    \item[\textbf{Problem 2.8.}] Confirm eq. (2.28) by applying the basic method to the mixer problem.

    \item[\textbf{Problem 2.9.}] Confirm eq. (2.30) by applying the rest of the Buckingham Pi theorem to the pendulum problem.

    \item[\textbf{Problem 2.10.}] Confirm eq. (2.32) by applying the rest of the Buckingham Pi theorem to the revised formulation of the pendulum problem.

    \item[\textbf{Problem 2.11.}] Apply the Buckingham Pi theorem to the revolution of two bodies in space, beginning with the functional equation (2.16).

    \item[\textbf{Problem 2.12.}] What happens when the Pi theorem is applied to the two-body problem, but beginning now with the functional equation (2.14)?

    \item[\textbf{Problem 2.13.}] Consider a string of length $l$ that connects a rock of mass $m$ to a fixed point while the rock whirls in a circle at speed $v$. Use the basic method of dimensional analysis to show that the tension $T$ in the string is determined by the dimensionless group:
    \begin{equation}
        \frac{TL}{m v^2} = \text{constant}
    \end{equation}

    \item[\textbf{Problem 2.14.}] Apply the Buckingham Pi theorem to confirm the analysis of Problem 2.13.

    \item[\textbf{Problem 2.15.}] The speed of sound in a gas, $c$, depends on the gas pressure $p$ and on the gas mass density $\rho$. Use dimensional analysis to determine how $c$, $p$, and $\rho$ are related.

    \item[\textbf{Problem 2.16.}] A dimensionless grouping called the Weber number, $We$, is used in fluid mechanics to relate a flowing fluid's surface tension, $\sigma$, which has dimensions of force/length, to the fluid's speed, $v$, density, $\rho$, and a characteristic length, $l$. Use dimensional analysis to find that number.

    \item[\textbf{Problem 2.17.}] A pendulum swings in a viscous fluid. How many groups are needed to relate the usual pendulum variables to the fluid viscosity, $\mu$, the fluid mass density, $\rho$, and the diameter $d$ of the pendulum? Find those groups.

    \item[\textbf{Problem 2.18.}] The volume flow rate $Q$ of fluid through a pipe is thought to depend on the pressure drop per unit length, $\Delta p / l$, the pipe diameter, $d$, and the fluid viscosity, $\mu$. Use the basic method of dimensional analysis to determine the relation:
    \begin{equation}
        Q = (\text{constant}) \left(\frac{d^4}{\mu}\right) \left(\frac{\Delta p}{l}\right)
    \end{equation}

    \item[\textbf{Problem 2.19.}] Apply the Buckingham Pi theorem to confirm the analysis of Problem 2.18.

    \item[\textbf{Problem 2.20.}] When flow in a pipe with a rough inner wall (perhaps due to a build-up of mineral deposits) is considered, several variables must be considered, including the fluid speed $v$, its density $\rho$ and viscosity $\mu$, and the pipe length $l$ and diameter $d$. The average variation $e$ of the pipe radius can be taken as a measure of the roughness of the pipe's inner surface. Determine the dimensionless groups needed to determine how the pressure drop $\Delta p$ depends on these variables.

    \item[\textbf{Problem 2.21.}] Use dimensional analysis to determine how the speed of sound in steel depends on the modulus of elasticity, $E$, and the mass density, $\rho$. (The modulus of elasticity of steel is, approximately and in British units, $30 \times 10^6$ psi.)

    \item[\textbf{Problem 2.22.}] The flexibility (the deflection per unit load) or compliance $C$ of a beam having a square cross-section $d \times d$ depends on the beam's length $l$, its height and width, and its material's modulus of elasticity $E$. Use the basic method of dimensional analysis to show that:
    \begin{equation}
        CEd = F_{CEd}\left(\frac{l}{d}\right)
    \end{equation}

    \item[\textbf{Problem 2.23.}] Experiments were conducted to determine the specific form of the function $F_{CEd}(l/d)$ found in Problem 2.22. In these experiments it was found that a plot of $\log_{10} (CEd)$ against $\log_{10} (l / d)$ has a slope of 3 and an intercept on the $\log_{10} (CEd)$ scale of $-0.60$. Show that the deflection under a load $P$ can be given in terms of the second moment of area $I$ of the cross section ($I = d^4/12$) as:
    \begin{equation}
        \text{deflection} = \text{load} \times \text{compliance} = P \times C = \frac{Pl^3}{48EI}
    \end{equation}

\end{enumerate}